\documentclass[11pt]{article}

\usepackage[a4paper,margin=1in]{geometry}
\usepackage{amsmath,amssymb}
\usepackage{booktabs}

\title{Econometrics -- Problem Set 2}
\author{Travis Bruffy}
\date{}

\begin{document}
\maketitle

\section*{Problem 1(a)}

Define $lwage_i=\log(wage_i)$. For each education level
$e\in\{12,\ldots,17\}$ and age $a\in\{26,\ldots,30\}$,
define
\[
d_{e,a,i}
=\mathbf{1}\{education_i=e,\ age_i=a\}.
\]

We estimate the regression
\[
lwage_i
=\sum_{e=12}^{17}\sum_{a=26}^{30}
\alpha_{e,a}d_{e,a,i}+\epsilon_i.
\]
No additional intercept is included because the
30 dummy variables sum to one for every observation.

Since $E(\epsilon_i\mid education_i,age_i)=0$,
\[
\alpha_{e,a}
=f(e,a)
=E[lwage_i\mid education_i=e,\ age_i=a].
\]
Each coefficient is therefore the population mean
log wage for its education--age group. Its OLS
estimate is the corresponding sample mean log wage.
This specification permits a separate conditional
mean for every group.

Table~\ref{tab:ps1a} reports the estimates and
conventional standard errors, which assume a
common error variance across groups.

\begin{table}[htbp]
\centering
\small
\caption{Education--age group estimates}
\label{tab:ps1a}
\begin{tabular}{lrrr}
\toprule
 & $\widehat{\alpha}$ & Standard error & Group size \\
\midrule
d1226 & 9.4513 & 0.0908 & 85 \\
d1227 & 9.3004 & 0.0850 & 97 \\
d1228 & 9.5379 & 0.0833 & 101 \\
d1229 & 9.3858 & 0.0854 & 96 \\
d1230 & 9.6283 & 0.0877 & 91 \\
d1326 & 9.6121 & 0.1920 & 19 \\
d1327 & 9.7559 & 0.1920 & 19 \\
d1328 & 9.3782 & 0.1745 & 23 \\
d1329 & 9.8200 & 0.2237 & 14 \\
d1330 & 9.4465 & 0.2030 & 17 \\
d1426 & 9.5807 & 0.2237 & 14 \\
d1427 & 9.7944 & 0.1708 & 24 \\
d1428 & 9.6422 & 0.2161 & 15 \\
d1429 & 9.7122 & 0.1745 & 23 \\
d1430 & 9.6454 & 0.1708 & 24 \\
d1526 & 9.5346 & 0.2416 & 12 \\
d1527 & 9.3856 & 0.2321 & 13 \\
d1528 & 9.5732 & 0.2647 & 10 \\
d1529 & 9.5404 & 0.2524 & 11 \\
d1530 & 9.9195 & 0.2416 & 12 \\
d1626 & 10.0705 & 0.1674 & 25 \\
d1627 & 9.9444 & 0.1871 & 20 \\
d1628 & 9.9920 & 0.1480 & 32 \\
d1629 & 10.1785 & 0.1307 & 41 \\
d1630 & 10.0706 & 0.1641 & 26 \\
d1726 & 9.9583 & 0.4832 & 3 \\
d1727 & 10.0458 & 0.2237 & 14 \\
d1728 & 10.1656 & 0.2959 & 8 \\
d1729 & 9.6176 & 0.4185 & 4 \\
d1730 & 10.0160 & 0.3743 & 5 \\
\bottomrule
\end{tabular}

\end{table}

\end{document}